\documentclass[a4paper,12pt]{article}
\usepackage[T2A]{fontenc}
\usepackage[utf8]{inputenc}
\usepackage[russian]{babel}
\usepackage{amsmath,amssymb,mathtools}
\usepackage{geometry}
\geometry{left=24mm,right=23mm,top=22mm,bottom=22mm}\setlength{\headheight}{14pt}
\usepackage{microtype}
\usepackage{array,booktabs,tabularx}
\usepackage{xcolor}
\usepackage{hyperref}
\usepackage{tikz}
\usepackage{fancyhdr}
\usepackage{enumitem}
\usepackage{parskip}
\usepackage{setspace}
\definecolor{accent}{HTML}{007B80}
\definecolor{deep}{HTML}{174C50}
\definecolor{muted}{HTML}{484848}
\definecolor{errorred}{HTML}{A32329}
\hypersetup{colorlinks=true,linkcolor=deep,pdfauthor={Лёвин Артём Александрович},pdftitle={Ревью домашней работы — Матвей Горбачев}}
\setstretch{1.16}
\setlist[itemize]{leftmargin=*,itemsep=3pt,topsep=4pt}
\pagestyle{fancy}
\fancyhf{}
\fancyhead[L]{\small\color{muted} Математика · оптимизация}
\fancyhead[R]{\small\color{muted} 17.11.2026}
\fancyfoot[C]{\small\thepage}
\renewcommand{\headrulewidth}{0.35pt}
\renewcommand{\footrulewidth}{0pt}
\newcommand{\markerror}{\textcolor{errorred}{\textbf{Ошибка:}}}
\newcommand{\keyidea}{\textcolor{deep}{\textbf{Ключевая идея:}}}
\newcommand{\outcome}{\textcolor{accent}{\textbf{Итог:}}}
\newcommand{\sublabel}[1]{\par\vspace{0.35em}\noindent\textbf{#1}\par\vspace{0.05em}}
\begin{document}
\thispagestyle{empty}
\begin{center}
\vspace*{23mm}
{\color{deep}\large МАТЕМАТИКА}\par
\vspace{22mm}
{\Huge\color{muted} Ревью домашней работы}\par
\vspace{10mm}
{\Large\color{accent} Оптимизация: распределение ресурсов\par и производственный план}\par
\vspace{16mm}
{\normalsize Задачи с ограничениями и целочисленным ответом}\par
\vfill
{\large Матвей Горбачев}\par
\vspace{11mm}
{17 ноября 2026 г.}\par
\vspace{17mm}
{\small Лёвин Артём Александрович}\par
{\small эксклюзивно для Матвея Горбачева}\par
\vspace{21mm}
\begin{tikzpicture}[x=1cm,y=1cm]
\draw[accent,line width=0.8pt] (0,0) .. controls (1.2,0.55) and (2.6,-0.40) .. (3.9,0)
 .. controls (5.1,0.45) and (6.4,-0.38) .. (7.7,0)
 .. controls (8.8,0.32) and (9.2,0.12) .. (9.8,0);
\end{tikzpicture}
\end{center}
\clearpage
\section*{Сводная таблица}
\addcontentsline{toc}{section}{Сводная таблица}
\noindent\small В таблицу включены задания, по которым на предоставленных фотографиях отсутствует решение. Для каждого задания ниже приведены объяснение метода, полный разбор и одно упражнение на закрепление.\normalsize\par
\vspace{8mm}
\renewcommand{\arraystretch}{1.45}
{\footnotesize\begin{tabularx}{\linewidth}{@{}p{9mm}p{39mm}X p{18mm}p{31mm}@{}}
\toprule
№ & Тема & Суть недочёта & Ваш ответ & Правильный ответ \\
\midrule
\hyperref[task:9]{9} & Максимальный выпуск & Решение на фотографиях не представлено & не указан & 26 изделий: 18 и 8 \\
\addlinespace
\hyperref[task:10]{10} & Производственный план & Решение на фотографиях не представлено & не указан & 50 изделий: 32 и 18 \\
\bottomrule
\end{tabularx}}\par
\vfill
\noindent\textcolor{deep}{\textbf{Основной приём.}} Если необходимо получить максимальный суммарный выпуск при ограничении на сумму квадратов, удобно сначала обозначить общее число изделий через $N$, затем выразить затраты через $N$ и одно из распределяемых количеств. Выделение полного квадрата даёт верхнюю оценку для $N$; остаётся проверить, можно ли достигнуть её при целых объёмах производства.
\clearpage
\thispagestyle{empty}
\vspace*{0.33\textheight}
\begin{center}
{\color{deep}\Huge Разбор ошибок и пропусков}\par
\vspace{5mm}
{\color{muted}\large Максимизация выпуска при бюджетных ограничениях}
\end{center}
\clearpage
\phantomsection\label{task:9}
\section*{Задание 9. Максимальный выпуск при ограниченном бюджете}
\sublabel{Текст задания}
Два завода производят одинаковые изделия. Если первый завод выпускает $t_1$ изделий за неделю, оплата труда составляет $160t_1^2$ рублей; если второй выпускает $t_2$ изделий, оплата труда составляет $360t_2^2$ рублей. Общий недельный бюджет на оплату труда --- не более $74\,880$ рублей. Найдите максимальный суммарный выпуск и распределение изделий по заводам. Числа $t_1,t_2$ должны быть целыми неотрицательными.

\sublabel{Ваше решение и ответ}
На четырёх переданных фотографиях нет записи решения задания №\,9. Ответ не указан.

\sublabel{Ошибка}
\markerror\ Пропущено решение задания. Проверить собственный математический ход и обнаружить первый неверный переход невозможно.

\sublabel{Где именно возникает проблема}
Последний правильный шаг и первый неправильный шаг выделить нельзя: записанных шагов нет. Требуется самостоятельно составить ограничение бюджета и доказать, что больший выпуск невозможен.

\sublabel{Почему недостаточно подобрать подходящие числа}
Даже если найти два количества изделий, при которых весь бюджет расходуется полностью, это ещё не доказывает максимальность суммы. Нужны \emph{две части}: верхняя оценка общего выпуска и пример, на котором эта оценка достигается.

\sublabel{Ключевая идея}
\keyidea\ Ввести $N=t_1+t_2$ и выделить полный квадрат в выражении общих затрат. Квадрат неотрицателен, поэтому получится оценка для $N$, не зависящая от конкретного распределения.

\sublabel{Исправление и полное правильное решение}
\textbf{Шаг 1.} Составим бюджетное ограничение:
\[
160t_1^2+360t_2^2\le 74\,880,\qquad
 t_1,t_2\in\mathbb Z_{\ge0}.
\]
Разделим обе части на положительное число $40$:
\[
4t_1^2+9t_2^2\le1872.
\]

\textbf{Шаг 2.} Обозначим суммарный выпуск через $N=t_1+t_2$. Тогда $t_2=N-t_1$, поэтому
\begin{align*}
4t_1^2+9t_2^2
 &=4t_1^2+9(N-t_1)^2\\
 &=4t_1^2+9N^2-18Nt_1+9t_1^2\\
 &=13t_1^2-18Nt_1+9N^2.
\end{align*}

\textbf{Шаг 3.} Выделим полный квадрат:
\begin{align*}
13t_1^2-18Nt_1+9N^2
 &=13\left(t_1-\frac{9N}{13}\right)^2
   +\frac{36N^2}{13}.
\end{align*}
Первое слагаемое неотрицательно при любых действительных значениях $t_1$ и $N$. Следовательно,
\[
4t_1^2+9t_2^2\ge\frac{36N^2}{13}.
\]
Совместим это с бюджетным ограничением:
\[
\frac{36N^2}{13}\le 1872
\quad\Longrightarrow\quad
N^2\le\frac{1872\cdot13}{36}=676.
\]
Так как $N\ge0$, получаем $N\le 26$. Значит, \emph{ни одно допустимое распределение не может дать более 26 изделий}.

\textbf{Шаг 4.} Проверим достижимость верхней оценки. Равенство в выделенном квадрате возможно, когда
\[
t_1=\frac{9N}{13}.
\]
При $N=26$ получаем $t_1=18$, $t_2=26-18=8$. Оба числа --- целые и неотрицательные.

\sublabel{Проверка}
\[
160\cdot18^2+360\cdot8^2
=160\cdot324+360\cdot64
=51\,840+23\,040=74\,880.
\]
Затраты равны разрешённому бюджету, а выпуск равен найденной верхней границе.

\sublabel{Итог}
\outcome\ Максимальный выпуск --- $\boxed{26}$ изделий в неделю: $\boxed{18}$ на первом заводе и $\boxed{8}$ на втором.

\sublabel{Задача на закрепление}
Два цеха производят один вид изделий. При выпуске $a$ изделий первым цехом затраты составляют $100a^2$ рублей, вторым при выпуске $b$ изделий --- $400b^2$ рублей. Бюджет не превышает $32\,000$ рублей. Найдите наибольшее возможное число изделий и распределение выпуска, если $a$ и $b$ --- целые неотрицательные числа.
\clearpage
\phantomsection\label{task:10}
\section*{Задание 10. Оптимизация производственного плана}
\sublabel{Текст задания}
Компания владеет двумя предприятиями, выпускающими один вид продукции. При производстве $t_1$ изделий на первом предприятии недельные расходы на оплату труда составляют $180t_1^2$ рублей, а при производстве $t_2$ изделий на втором --- $320t_2^2$ рублей. Общий недельный бюджет ограничен суммой $288\,000$ рублей. Найдите максимально возможный суммарный выпуск и его распределение между предприятиями, если $t_1,t_2$ --- целые неотрицательные числа.

\sublabel{Ваше решение и ответ}
Решение задания №\,10 на предоставленных фотографиях отсутствует. Ответ не указан.

\sublabel{Ошибка}
\markerror\ Задание оставлено без решения. Достоверного ошибочного перехода в записи нет.

\sublabel{Где именно возникает проблема}
Не построена проверяемая математическая модель ограничения бюджета и не получено доказательство максимально возможного суммарного выпуска. Последний правильный и первый неправильный шаг указать нельзя из-за отсутствия записи.

\sublabel{Почему важно доказать максимальность}
Подставить несколько пар целых чисел и выбрать лучшую среди них недостаточно. Необходимо получить границу, выше которой суммарный выпуск в принципе не может подняться при данном бюджете, а затем построить допустимую пару, достигающую границы.

\sublabel{Ключевая идея}
\keyidea\ Применить тот же принцип, что в задаче №\,9, но с другими коэффициентами затрат. Выражение после раскрытия скобок снова превращается в сумму неотрицательного квадрата и величины, зависящей только от $N=t_1+t_2$.

\sublabel{Исправление и полное правильное решение}
\textbf{Шаг 1.} Запишем условие на расходы и допустимые значения:
\[
180t_1^2+320t_2^2\le 288\,000,
\qquad t_1,t_2\in\mathbb Z_{\ge0}.
\]
Разделим обе части на $20$:
\[
9t_1^2+16t_2^2\le14\,400.
\]

\textbf{Шаг 2.} Введём $N=t_1+t_2$ и заменим $t_2=N-t_1$:
\begin{align*}
9t_1^2+16t_2^2
 &=9t_1^2+16(N-t_1)^2\\
 &=9t_1^2+16N^2-32Nt_1+16t_1^2\\
 &=25t_1^2-32Nt_1+16N^2.
\end{align*}

\textbf{Шаг 3.} Выделим полный квадрат:
\[
25t_1^2-32Nt_1+16N^2
=25\left(t_1-\frac{16N}{25}\right)^2
 +\frac{144N^2}{25}.
\]
Так как квадрат неотрицателен,
\[
9t_1^2+16t_2^2\ge\frac{144N^2}{25}.
\]
Совместно с бюджетным ограничением это даёт
\[
\frac{144N^2}{25}\le14\,400
\quad\Longrightarrow\quad
N^2\le\frac{14\,400\cdot25}{144}=2500.
\]
Учитывая $N\ge0$, заключаем: $N\le50$.

\textbf{Шаг 4.} При $N=50$ равенство в выделенном квадрате требует
\[
t_1=\frac{16\cdot50}{25}=32,
\qquad t_2=50-32=18.
\]
Получены целые неотрицательные значения, то есть верхняя оценка достижима и при условии целочисленности.

\sublabel{Проверка}
\[
180\cdot32^2+320\cdot18^2
=180\cdot1024+320\cdot324
=184\,320+103\,680=288\,000.
\]
Ограничение выполнено с равенством. Суммарный выпуск $32+18=50$.

\sublabel{Итог}
\outcome\ Максимально возможный выпуск --- $\boxed{50}$ изделий в неделю: $\boxed{32}$ на первом предприятии и $\boxed{18}$ на втором.

\sublabel{Задача на закрепление}
Две мастерские выпускают одинаковые изделия. Недельные затраты при выпуске $u$ изделий первой мастерской равны $200u^2$ рублей, при выпуске $v$ изделий второй --- $450v^2$ рублей. Общий бюджет не превышает $93\,600$ рублей. Найдите максимальный общий выпуск и распределение заказов между мастерскими. Значения $u$ и $v$ --- целые неотрицательные числа.

\end{document}